
\section{Постановка на задачата}
Целта на упражнението е симулационно изследване на организацията и системните характеристики на базови конвейерни архитектури в средата GPSS World. Задачите обхващат:
\begin{enumerate}
    \item Разучаване на организацията и диспечерирането на конвейерни структури.
    \item Симулационно изследване на \textbf{линеен едно-функционален конвейер (ЕФК)} с постоянен такт при вариращ входен поток.
    \item Разработване на модел за \textbf{ЕФК с обратна връзка} при рекурсивна обработка.
    \item Разработване на модел за \textbf{многофункционален конвейер (МФК)} с два типа операции.
\end{enumerate}

\section{Резултати от експериментите}

\subsection{Задача 2: Линеен ЕФК с постоянен такт}

Изследван е линеен едно-функционален конвейер с $k = 4$ фази (функционални устройства PROC1--PROC4) и постоянен такт $t = 5$ моделни единици. Входният поток е с експоненциално разпределение. Проведени са 5 симулационни експеримента с вариращо средно време между постъпванията (IAT), всеки симулира 300 завършени задачи. Теоретичната натовареност на всяка фаза е $\rho = t / \mathrm{IAT} = 5 / \mathrm{IAT}$.

\textbf{Код на симулационния модел (GPSS, пример при IAT\,=\,7):}
\begin{lstlisting}[language=bash]
EXPON FUNCTION RN1,C24
0,0/.1,.104/.2,.222/.3,.355/.4,.509/.5,.69/.6,.915/.7,1.2
.75,1.38/.8,1.6/.84,1.83/.88,2.12/.9,2.3/.92,2.52/.94,2.81
.95,2.99/.96,3.2/.97,3.5/.98,3.9/.99,4.6/.995,5.3/.998,6.2
.999,7/.9998,8

         GENERATE 7,FN$EXPON    ; IAT varied: 2, 4, 5, 7, 10
         QUEUE    WAIT1          ; Phase 1 queue
         SEIZE    PROC1          ; Occupy FU1
         DEPART   WAIT1
         ADVANCE  5              ; Fixed clock t = 5
         RELEASE  PROC1
         QUEUE    WAIT2          ; Phase 2 queue
         SEIZE    PROC2
         DEPART   WAIT2
         ADVANCE  5
         RELEASE  PROC2
         QUEUE    WAIT3          ; Phase 3 queue
         SEIZE    PROC3
         DEPART   WAIT3
         ADVANCE  5
         RELEASE  PROC3
         QUEUE    WAIT4          ; Phase 4 queue
         SEIZE    PROC4
         DEPART   WAIT4
         ADVANCE  5
         RELEASE  PROC4
         TERMINATE 1
         START    300
\end{lstlisting}

\subsubsection{Резултати}

\begin{table}[H]
    \centering
    \caption{Натовареност на ФУ при вариращ входен поток (300 задачи)}
    \begin{tabular}{|c|c|c|c|c|c|c|}
        \hline
        \textbf{IAT} & \textbf{$\rho_\mathrm{теор}$} & \textbf{PROC1} & \textbf{PROC2} & \textbf{PROC3} & \textbf{PROC4} & \textbf{Ср. чакане [м.е.]} \\
        \hline
        2  & $>$100\% & 99{,}9\% & 99{,}5\% & 99{,}2\% & 98{,}9\% & 434{,}0 \\
        \hline
        4  & $>$100\% & 99{,}0\% & 98{,}6\% & 98{,}3\% & 98{,}0\% & 149{,}5 \\
        \hline
        5  & 100\%    & 93{,}4\% & 93{,}2\% & 93{,}0\% & 92{,}7\% & 17{,}5  \\
        \hline
        7  & 71{,}4\% & 70{,}6\% & 70{,}6\% & 70{,}4\% & 70{,}2\% & 5{,}5   \\
        \hline
        10 & 50{,}0\% & 51{,}3\% & 51{,}2\% & 51{,}0\% & 50{,}9\% & 2{,}4   \\
        \hline
    \end{tabular}
\end{table}

\begin{figure}[H]
    \centering
    \begin{tikzpicture}
        \begin{axis}[
            title={Натовареност на PROC1 спрямо средно IAT},
            xlabel={Средно IAT [мод. ед.]},
            ylabel={Натовареност (\%)},
            grid=major,
            width=0.8\linewidth,
            height=6cm,
            xmin=1, xmax=11,
            ymin=0, ymax=110,
            xtick={2,4,5,7,10},
            legend pos=south east,
        ]
        \addplot[color=blue, mark=*, thick]
            coordinates {(2,99.9)(4,99.0)(5,93.4)(7,70.6)(10,51.3)};
        \addplot[color=red, dashed, thick, mark=square*]
            coordinates {(2,100)(4,100)(5,100)(7,71.4)(10,50.0)};
        \legend{Симулация (PROC1), Теория}
        \end{axis}
    \end{tikzpicture}
    \caption{Зависимост на натовареността от средното IAT}
    \label{fig:task2_util}
\end{figure}

При $\mathrm{IAT} \leq t = 5$ конвейерът е претоварен ($\rho \geq 1$) и опашката пред PROC1 нараства неограничено. При $\mathrm{IAT} = 2$ максималният брой задачи в WAIT1 достига 446, а средното чакане е 434{,}0 мод. ед. При $\mathrm{IAT} > t$ конвейерът работи в стационарен режим и натовареността съответства точно на теоретичната.

\subsection{Задача 3: ЕФК с обратна връзка}

Изследван е три-фазен ЕФК ($k = 3$, такт $t = 5$) с вероятност за обратна връзка $p = 0{,}4$, имитираща рекурсивни изчисления. Задача, преминала всички три фази, с вероятност 40\% повторно постъпва в конвейера от начало. Средното IAT на входния поток е 7 мод. ед.

Ефективната натовареност на всяка фаза се определя по:
\begin{equation}
    \rho_\mathrm{eff} = \frac{t}{\mathrm{IAT} \cdot (1-p)} = \frac{5}{7 \cdot 0{,}6} \approx 1{,}19 > 1
\end{equation}

Тъй като $\rho_\mathrm{eff} > 1$, системата е в режим на претоварване.

\textbf{Код на симулационния модел (GPSS):}
\begin{lstlisting}[language=bash]
EXPON FUNCTION RN1,C24
0,0/.1,.104/.2,.222/.3,.355/.4,.509/.5,.69/.6,.915/.7,1.2
.75,1.38/.8,1.6/.84,1.83/.88,2.12/.9,2.3/.92,2.52/.94,2.81
.95,2.99/.96,3.2/.97,3.5/.98,3.9/.99,4.6/.995,5.3/.998,6.2
.999,7/.9998,8

         GENERATE 7,FN$EXPON    ; Inter-arrival mean = 7

PIPEIN   QUEUE    WAIT1          ; Phase 1 queue (feedback entry point)
         SEIZE    PROC1
         DEPART   WAIT1
         ADVANCE  5              ; Fixed clock t = 5
         RELEASE  PROC1
         QUEUE    WAIT2          ; Phase 2 queue
         SEIZE    PROC2
         DEPART   WAIT2
         ADVANCE  5
         RELEASE  PROC2
         QUEUE    WAIT3          ; Phase 3 queue
         SEIZE    PROC3
         DEPART   WAIT3
         ADVANCE  5
         RELEASE  PROC3
         TRANSFER .4,,PIPEIN    ; 40% feedback -> re-enters pipeline
         TERMINATE 1             ; 60% exit
         START    300
\end{lstlisting}

\subsubsection{Резултати}

\begin{table}[H]
    \centering
    \caption{Резултати от симулация на ЕФК с обратна връзка ($p = 0{,}4$, IAT\,=\,7)}
    \begin{tabular}{|l|c|}
        \hline
        \textbf{Показател} & \textbf{Стойност} \\
        \hline
        Генерирани задачи & 391 \\
        \hline
        Общо входове в конвейера & 639 \\
        \hline
        Средно прохода на задача & 1{,}63 \quad (теория: $1/(1-0{,}4) = 1{,}67$) \\
        \hline
        Натовареност PROC1 & 98{,}7\% \\
        \hline
        Натовареност PROC2 & 98{,}5\% \\
        \hline
        Натовареност PROC3 & 98{,}4\% \\
        \hline
        Максимална опашка WAIT1 & 90 задачи \\
        \hline
        Среден брой задачи в WAIT1 & 42{,}6 \\
        \hline
        Средно чакане WAIT1 & 185{,}6 мод. ед. \\
        \hline
    \end{tabular}
\end{table}

Средният брой проходи на задача (1{,}63) е в отлично съответствие с теоретичната стойност 1{,}67. Обратната връзка драстично увеличава натоварването на конвейера и предизвиква критично нарастване на опашката пред PROC1 (максимум 90 задачи, средно чакане 185{,}6 мод. ед.).

\subsection{Задача 4: Многофункционален конвейер (МФК)}

Изследван е динамичен МФК с четири функционални устройства и два независими входни потока:
\begin{itemize}
    \item \textbf{Операция A (целочислено събиране):} PROC1 $\to$ PROC2 $\to$ PROC3, IAT\,=\,8
    \item \textbf{Операция B (умножение с плаваща запетая):} PROC1 $\to$ PROC3 $\to$ PROC4, IAT\,=\,12
\end{itemize}

Такт на всяко ФУ: $t = 5$ мод. ед. Устройствата PROC1 и PROC3 са \textbf{споделени} между двете операции, моделирайки колизионен хазарт в МФК.

\textbf{Код на симулационния модел (GPSS):}
\begin{lstlisting}[language=bash]
EXPON FUNCTION RN1,C24
0,0/.1,.104/.2,.222/.3,.355/.4,.509/.5,.69/.6,.915/.7,1.2
.75,1.38/.8,1.6/.84,1.83/.88,2.12/.9,2.3/.92,2.52/.94,2.81
.95,2.99/.96,3.2/.97,3.5/.98,3.9/.99,4.6/.995,5.3/.998,6.2
.999,7/.9998,8

; --- Op A: Integer ADD (PROC1 -> PROC2 -> PROC3), IAT = 8 ---
         GENERATE 8,FN$EXPON
         QUEUE    WA1
         SEIZE    PROC1
         DEPART   WA1
         ADVANCE  5
         RELEASE  PROC1
         QUEUE    WA2
         SEIZE    PROC2
         DEPART   WA2
         ADVANCE  5
         RELEASE  PROC2
         QUEUE    WA3
         SEIZE    PROC3
         DEPART   WA3
         ADVANCE  5
         RELEASE  PROC3
         TERMINATE 1

; --- Op B: Float MULT (PROC1 -> PROC3 -> PROC4), IAT = 12 ---
         GENERATE 12,FN$EXPON
         QUEUE    WB1
         SEIZE    PROC1
         DEPART   WB1
         ADVANCE  5
         RELEASE  PROC1
         QUEUE    WB3
         SEIZE    PROC3
         DEPART   WB3
         ADVANCE  5
         RELEASE  PROC3
         QUEUE    WB4
         SEIZE    PROC4
         DEPART   WB4
         ADVANCE  5
         RELEASE  PROC4
         TERMINATE 1
         START    400
\end{lstlisting}

\subsubsection{Резултати}

\begin{table}[H]
    \centering
    \caption{Натовареност на ФУ в МФК (400 задачи)}
    \begin{tabular}{|l|c|c|c|}
        \hline
        \textbf{ФУ} & \textbf{Натовареност} & \textbf{Тип} & \textbf{$\rho_\mathrm{теор}$} \\
        \hline
        PROC1 & 96{,}8\% & Споделен (Op A + Op B) & $5/8 + 5/12 \approx 104\%$ \\
        \hline
        PROC2 & 61{,}1\% & Само Op A              & $5/8 = 62{,}5\%$           \\
        \hline
        PROC3 & 96{,}3\% & Споделен (Op A + Op B) & $5/8 + 5/12 \approx 104\%$ \\
        \hline
        PROC4 & 35{,}1\% & Само Op B              & $5/12 \approx 41{,}7\%$    \\
        \hline
    \end{tabular}
\end{table}

\begin{figure}[H]
    \centering
    \begin{tikzpicture}
        \begin{axis}[
            ybar,
            title={Натовареност на ФУ в МФК},
            ylabel={Натовареност (\%)},
            symbolic x coords={PROC1, PROC2, PROC3, PROC4},
            xtick=data,
            ymin=0, ymax=110,
            nodes near coords,
            nodes near coords align={vertical},
            width=0.8\linewidth,
            height=6cm,
            bar width=1.2cm,
            enlarge x limits=0.25,
        ]
        \addplot[fill=blue!60] coordinates {
            (PROC1,96.8)(PROC2,61.1)(PROC3,96.3)(PROC4,35.1)
        };
        \end{axis}
    \end{tikzpicture}
    \caption{Натовареност на функционалните устройства в МФК}
    \label{fig:task4_util}
\end{figure}

Споделените устройства PROC1 и PROC3 са силно натоварени (над 96\%), докато несподелените PROC2 и PROC4 работят при значително по-ниска натовареност. Средното чакане пред PROC1 е 34{,}4 мод. ед. за Операция A и 36{,}0 мод. ед. за Операция B (максимални опашки 14 и 10 задачи съответно). PROC4 е с по-ниска натовареност от теоретичната поради закъсненията от тясното място при PROC3.

\section{Анализ и заключение}

Проведените симулационни експерименти демонстрират основните характеристики на конвейерните архитектури:
\begin{itemize}
    \item При \textbf{Задача 2} натовареността на всяка фаза на линейния ЕФК е еднаква и съответства точно на теоретичната стойност $\rho = t / \mathrm{IAT}$. При $\mathrm{IAT} \leq t$ системата влиза в режим на претоварване --- опашките нарастват неограничено, а средното чакане нараства рязко (от 2{,}4 мод. ед. при IAT\,=\,10 до 434{,}0 мод. ед. при IAT\,=\,2).

    \item При \textbf{Задача 3} обратната връзка с $p = 0{,}4$ увеличава ефективното натоварване до $\rho_\mathrm{eff} \approx 1{,}19$, което води до критично претоварване на конвейера. Средният брой проходи на задача (1{,}63) съответства отлично на теоретичната стойност $1/(1-p) = 1{,}67$, потвърждавайки коректността на модела.

    \item При \textbf{Задача 4} споделянето на PROC1 и PROC3 между двата потока води до тяхното пресищане (около 97\%), демонстрирайки колизионния хазарт в динамичен МФК. Несподелените устройства PROC2 и PROC4 работят при значително по-ниско натоварване. Резултатите потвърждават, че споделените ресурси са определящото тясно място в МФК архитектурата.
\end{itemize}
